\documentclass[../../main.tex]{subfiles}
\graphicspath{{\subfix{../../image/}}}

\begin{document}

\section*{习题三}\label{section:ch03-习题三}

\begin{exercise}\label{exercise:ch03-01}
直线上有一质点，每经一个单位时间，它分别以概率 $p$ 及 $q$ 向右或向左移动一格，若该质点在时刻 $0$ 从原点出发，而且每次移动是相互独立的，试用随机变量来描述这质点的运动（以 $S_n$ 表示时刻 $n$ 时质点的位置）。
\end{exercise}

\begin{exercise}\label{exercise:ch03-02}
设 $\xi$ 为 Bernoulli 试验中第一个游程（连续的成功或失败）的长，试求 $\xi$ 的概率分布。
\end{exercise}

\begin{exercise}\label{exercise:ch03-03}
$C$ 应取何值才能使下列数列成为概率分布：
\begin{enumerate}
\item $p_k=\dfrac{C}{N},\quad k=1,2,\cdots,N$；
\item $p_k=C\dfrac{\lambda^{k}}{k!},\quad k=1,2,\cdots,\lambda>0$。
\end{enumerate}
\end{exercise}

\begin{exercise}\label{exercise:ch03-04}
若分布函数定义为 $F(x)=P\{\xi\leqslant x\}$，试证这时的 $F(x)$ 具有下列性质：(i) 非降；(ii) $F(-\infty)=0$，$F(+\infty)=1$；(iii) 右连续。
\end{exercise}

\begin{exercise}\label{exercise:ch03-05}
若 $\zeta\sim N(0,1)$，试求常数 $a,b,c$ 使 (1) $a=P\{\zeta\geqslant 1.645\}$；(2) $P\{|\zeta|<b\}=95\%$；(3) $P\{|\zeta-c|>c\}=0.51$。
\end{exercise}

\begin{exercise}\label{exercise:ch03-06}
妊娠天数 $\xi$ 的分布函数为 $N(270,100)$，求 $\xi$ 落在下列范围的概率：(1) $(260,280)$；(2) 短于 $250$ 天；(3) 长于 $300$ 天。
\end{exercise}

\begin{exercise}\label{exercise:ch03-07}
若 $\xi$ 的分布函数为 $N(60,9)$，求分点 $x_1,x_2,x_3,x_4$，使 $\xi$ 落在 $(-\infty,x_1)$，$(x_1,x_2)$，$(x_2,x_3)$，$(x_3,x_4)$，$(x_4,\infty)$ 中的概率之比为 $7:24:38:24:7$。
\end{exercise}

\begin{exercise}\label{exercise:ch03-08}
在 Pascal 分布 $\dbinom{k-1}{r-1}p^{r}q^{k-r}$ 中，令 $k\Delta t=t$，$p=\lambda\Delta t$，试证当 $\Delta t\to 0$ 时，它能用 $\dfrac{\lambda(\lambda t)^{r-1}}{(r-1)!}\mathrm{e}^{-\lambda t}\cdot\Delta t$ 来逼近。（这可以解释为第 $r$ 次成功发生在 $(t,t+\Delta t)$ 中的概率，其密度函数正好是参数为 $r$ 的 Erlang 分布。用这种方法可以把课文中的对比严格化。）
\end{exercise}

\begin{exercise}\label{exercise:ch03-09}
在生存分析中，作为研究对象的是非负随机变量，它们的分布称为寿命分布。若 $\xi$ 是非负随机变量，其分布函数为 $F(x)$，密度函数为 $f(x)$，这时通常还引入\textbf{生存函数} $S(x)=P\{\xi\geqslant x\}$ 及\textbf{失效率函数} $\lambda(x)=\dfrac{f(x)}{1-F(x)}$。试导出 $S(x)$，$\lambda(x)$，$F(x)$ 及 $f(x)$ 之间的关系式，并以指数分布验证之。
\end{exercise}

\begin{exercise}\label{exercise:ch03-10}
设随机变量 $\xi$ 取值于 $[0,1]$，若 $P\{x\leqslant\xi<y\}$ 只与长度 $y-x$ 有关（对一切 $0\leqslant x\leqslant y\leqslant 1$），试证 $\xi$ 服从 $[0,1]$ 均匀分布。
\end{exercise}

\begin{exercise}\label{exercise:ch03-11}
若存在 $\varTheta$ 上的实值函数 $Q(\theta)$ 及 $D(\theta)$ 以及 $T(x)$ 及 $S(x)$，使
\begin{align*}f_\theta(x)=\exp\{Q(\theta)T(x)+D(\theta)+S(x)\},\end{align*}
则称 $\{f_\theta,\theta\in\varTheta\}$ 是一个单参数的指数族。证明 (1) 正态分布 $N(m_0,\sigma^{2})$，已知 $m_0$，关于参数 $\sigma$；(2) 正态分布 $N(m,\sigma_0^{2})$，已知 $\sigma_0$，关于参数 $m$；(3) Poisson 分布 $P(\lambda)$ 关于 $\lambda$ 是一个单参数的指数族。

但是 $[0,\theta]$ 上均匀分布，关于 $\theta$ 不是一个单参数的指数族。
\end{exercise}

\begin{exercise}\label{exercise:ch03-12}
定义二元函数
\begin{align*}F(x,y)=
\begin{cases}
1, & x+y>0,\\
0, & x+y\leqslant 0.
\end{cases}\end{align*}
验证此函数对每个变元非降，左连续，且满足分布函数性质 (ii)，但无法使\eqref{align:ch03-二元分布非负性}保持非负。
\end{exercise}

\begin{exercise}\label{exercise:ch03-13}
若 $f_1(x),f_2(y)$ 为分布密度，求为使 $f(x,y)=f_1(x)\times f_2(y)+h(x,y)$ 成为密度函数，$h(x,y)$ 必须而且只须满足什么条件。
\end{exercise}

\begin{exercise}\label{exercise:ch03-14}
若 $f_1(x),f_2(x),f_3(x)$ 是对应于分布函数 $F_1(x),F_2(x),F_3(x)$ 的密度函数，证明对于一切 $\alpha$（$-1<\alpha<1$），下列函数是密度函数，且具有相同的边际密度函数 $f_1(x),f_2(x),f_3(x)$：
\begin{align*}
f_\alpha(x_1,x_2,x_3)
&=f_1(x_1)f_2(x_2)f_3(x_3)\{1+\alpha[2F_1(x_1)-1]\\
&\quad\cdot[2F_2(x_2)-1][2F_3(x_3)-1]\}.
\end{align*}
\end{exercise}

\begin{exercise}\label{exercise:ch03-15}
若 $(\xi,\eta)$ 的联合概率分布为
\begin{align*}\begin{array}{c|ccc}
\eta\backslash\xi & -1 & 0 & 1\\
\hline
-1 & a & 0 & 0.2\\
0 & 0.1 & b & 0.1\\
1 & 0 & 0.2 & c
\end{array}\end{align*}
且 $P\{\xi\eta\neq 0\}=0.4$，$P\{\eta\leqslant 0\mid\xi\leqslant 0\}=\dfrac{2}{3}$，试求：
\begin{enumerate}
\item $a,b,c$ 之值；
\item $\xi$ 及 $\eta$ 的边际概率分布；
\item $\xi+\eta$ 的概率分布。
\end{enumerate}
\end{exercise}

\begin{exercise}\label{exercise:ch03-16}
若 $(\xi,\eta)$ 的密度函数为
\begin{align*}p(x,y)=
\begin{cases}
A\mathrm{e}^{-(2x+y)}, & x>0,\ y>0,\\
0, & \text{其他}.
\end{cases}\end{align*}
试求：(1) 常数 $A$；(2) $P\{\xi<2,\eta<1\}$；(3) $\xi$ 的边际分布函数；(4) $P\{\xi+\eta<2\}$；(5) $p(x\mid y)$；(6) $P\{\xi<2\mid\eta<1\}$。
\end{exercise}

\begin{exercise}\label{exercise:ch03-17}
若 $P\{\mu=m,\nu=n\}=\dfrac{(\lambda p)^{m}(\lambda-\lambda p)^{n-m}}{m!(n-m)!}\mathrm{e}^{-\lambda}$，$m=0,1,2,\cdots,n$；$n=0,1,2,\cdots$，试求：(1) $P\{\nu=n\}$；(2) $P\{\mu=m\}$；(3) $P\{\mu=m\mid\nu=n\}$；(4) $P\{\nu-\mu=k\}$。
\end{exercise}

\begin{exercise}\label{exercise:ch03-18}
设二维随机变量 $(\xi,\eta)$ 的联合密度为
\begin{align*}p(x,y)=\frac{1}{\Gamma(k_1)\Gamma(k_2)}x^{k_1-1}(y-x)^{k_2-1}\mathrm{e}^{-y},\end{align*}
$k_1>0,k_2>0,0<x\leqslant y<\infty$。试求 $\xi$ 与 $\eta$ 的边际分布密度。
\end{exercise}

\begin{exercise}\label{exercise:ch03-19}
试证 $p(x,y)=K\mathrm{e}^{-(ax^{2}+2bxy+cy^{2})}$ 为密度函数的充要条件为 $a>0,c>0,b^{2}-ac<0$，$K=\dfrac{1}{\pi}\sqrt{ac-b^{2}}$。
\end{exercise}

\begin{exercise}\label{exercise:ch03-20}
\begin{enumerate}
\item 若 $(\xi,\eta)$ 的联合密度函数为
\begin{align*}f(x,y)=
\begin{cases}
4xy, & 0\leqslant x\leqslant 1,\ 0\leqslant y\leqslant 1,\\
0, & \text{其他},
\end{cases}\end{align*}
问 $\xi$ 与 $\eta$ 是否相互独立？
\item 若 $(\xi,\eta)$ 的联合密度函数为
\begin{align*}g(x,y)=
\begin{cases}
8xy, & 0\leqslant x\leqslant y,\ 0\leqslant y\leqslant 1,\\
0, & \text{其他},
\end{cases}\end{align*}
问 $\xi$ 与 $\eta$ 是否相互独立？
\end{enumerate}
\end{exercise}

\begin{exercise}\label{exercise:ch03-21}
若 $\xi,\eta$ 相互独立且皆以概率 $\dfrac{1}{2}$ 取值 $+1$ 及 $-1$，令 $\zeta=\xi\eta$，试证 $\xi,\eta,\zeta$ 两两独立但不相互独立。
\end{exercise}

\begin{exercise}\label{exercise:ch03-22}
设 $(\xi,\eta)$ 具有联合密度函数
\begin{align*}p(x,y)=
\begin{cases}
\dfrac{1+xy}{4}, & |x|<1,\ |y|<1,\\
0, & \text{其他}.
\end{cases}\end{align*}
试证 $\xi$ 与 $\eta$ 不独立，但 $\xi^{2}$ 与 $\eta^{2}$ 是相互独立的。
\end{exercise}

\begin{exercise}\label{exercise:ch03-23}
若每次试验中出现 $A_1,A_2,A_3$ 的概率分别为 $p_1,p_2$ 及 $p_3$，而且 $p_1+p_2+p_3=1$，共进行 $n$ 次独立试验，以 $\mu_1,\mu_2,\mu_3$ 分别记 $A_1,A_2,A_3$ 出现的次数，试求：
\begin{enumerate}
\item $(\mu_1,\mu_2)$ 的联合概率分布；
\item $\mu_1$ 的概率分布；
\item $P\{\mu_2=k_2\mid\mu_1=k_1\}$。
\end{enumerate}
\end{exercise}

\begin{exercise}\label{exercise:ch03-24}
袋中装 $i$ 号球 $N_i$ 只，$i=1,2,3$，$N_1+N_2+N_3=N$，从中随机摸出 $n$ 只，若以 $\mu_1,\mu_2,\mu_3$ 分别记 $1,2,3$ 号球出现的次数，试求：
\begin{enumerate}
\item $(\mu_1,\mu_2)$ 的联合概率分布；
\item $\mu_1$ 的概率分布；
\item $P\{\mu_2=n_2\mid\mu_1=n_1\}$。
\end{enumerate}
\end{exercise}

\begin{exercise}\label{exercise:ch03-25}
若 $\xi_1$ 与 $\xi_2$ 是独立随机变量，且 $\xi_1\sim B(n_1,p)$，$\xi_2\sim B(n_2,p)$，试直接证明：
\begin{enumerate}
\item $\xi_1+\xi_2\sim B(n_1+n_2,p)$；
\item $P\{\xi_1=k\mid\xi_1+\xi_2=n\}=\dfrac{\dbinom{n_1}{k}\dbinom{n_2}{n-k}}{\dbinom{n_1+n_2}{n}}$。
\end{enumerate}
\end{exercise}

\begin{exercise}\label{exercise:ch03-26}
若 $\xi_1$ 与 $\xi_2$ 是独立随机变量，均服从 Poisson 分布，参数分别为 $\lambda_1$ 及 $\lambda_2$，试直接证明：
\begin{enumerate}
\item $\xi_1+\xi_2$ 具有 Poisson 分布，参数为 $\lambda_1+\lambda_2$；
\item $P\{\xi_1=k\mid\xi_1+\xi_2=n\}=\dbinom{n}{k}\left(\dfrac{\lambda_1}{\lambda_1+\lambda_2}\right)^{k}\left(\dfrac{\lambda_2}{\lambda_1+\lambda_2}\right)^{n-k}$。
\end{enumerate}
\end{exercise}

\begin{exercise}\label{exercise:ch03-27}
设 $\xi$ 的密度函数为 $p(x)$，求下列随机变量的分布密度函数：(1) $\eta=\xi^{-1}$；(2) $\eta=\operatorname{tg}\xi$；(3) $\eta=|\xi|$。
\end{exercise}

\begin{exercise}\label{exercise:ch03-28}
设 $\xi$ 与 $\eta$ 相互独立且服从同一几何分布，令 $\zeta=\max(\xi,\eta)$，试求 (1) $(\zeta,\xi)$ 的联合概率分布；(2) $\zeta$ 的概率分布；(3) $\xi$ 关于 $\zeta$ 的条件概率分布。
\end{exercise}

\begin{exercise}\label{exercise:ch03-29}
若 $\xi,\eta$ 为相互独立的分别服从 $[0,1]$ 均匀分布的随机变量，试求 $\zeta=\xi+\eta$ 的分布密度函数。
\end{exercise}

\begin{exercise}\label{exercise:ch03-30}
在 $(0,a)$ 线段上随机投掷两点，试求两点间距离的分布函数。
\end{exercise}

\begin{exercise}\label{exercise:ch03-31}
若 $\xi$ 与 $\eta$ 相互独立，均服从参数为 $\dfrac{1}{\alpha}$ 的指数分布 $\operatorname{Exp}\!\left(\dfrac{1}{\alpha}\right)$，试证 $\zeta=\xi-\eta$ 服从 Laplace 分布：
\begin{align*}p(x)=\frac{1}{2\alpha}\mathrm{e}^{-\frac{|x|}{\alpha}},\qquad -\infty<x<\infty.\end{align*}
\end{exercise}

\begin{exercise}\label{exercise:ch03-32}
若 $\xi$ 与 $\eta$ 相互独立，分别服从 $N(0,1)$，试证 $\psi=\dfrac{\xi}{\eta}$ 服从 Cauchy 分布。
\end{exercise}

\begin{exercise}\label{exercise:ch03-33}
若 $\xi_1,\xi_2,\cdots,\xi_n$ 相互独立，且皆服从指数分布，参数分别为 $\lambda_1,\lambda_2,\cdots,\lambda_n$，试求 $\eta=\min(\xi_1,\xi_2,\cdots,\xi_n)$ 的分布。
\end{exercise}

\begin{exercise}\label{exercise:ch03-34}
通称下列分布函数为 Weibull 分布：
\begin{align*}F(x)=
\begin{cases}
1-\mathrm{e}^{-\lambda x^{\alpha}}, & x>0,\\
0, & x\leqslant 0.
\end{cases}\end{align*}
这是 Weibull 在研究金属材料的疲劳寿命中导出的，在可靠性研究中有广泛应用。

若 $\xi_1,\xi_2,\cdots,\xi_n$ 相互独立，相同分布，并以 $\xi_1^{*}$ 记它们的最小值，(1) 当 $\xi_i\sim F(x)$ 时，试求 $\xi_1^{*}$ 的分布函数。(2) 若 $\xi_1^{*}\sim F(x)$，试导出 $\xi_i$ 的分布函数。

链条的寿命取决于最弱环节，试说明上述概率结论的实际含意。
\end{exercise}

\begin{exercise}\label{exercise:ch03-35}
若 $(\xi,\eta)$ 服从二元正态分布，参数 $\mu_1,\mu_2,\sigma_1^{2},\sigma_2^{2},\rho$，以 $D(\lambda)$ 记下面椭圆的内部：
\begin{align*}\frac{(x-\mu_1)^{2}}{\sigma_1^{2}}-\frac{2\rho(x-\mu_1)(y-\mu_2)}{\sigma_1\sigma_2}+\frac{(y-\mu_2)^{2}}{\sigma_2^{2}}=\lambda^{2}.\end{align*}
试求 $P\{(\xi,\eta)\in D(\lambda)\}$。
\end{exercise}

\begin{exercise}\label{exercise:ch03-36}
若气体分子的速度是随机向量 $V=(X,Y,Z)$，各分量相互独立，且均服从 $N(0,\sigma^{2})$，试证 $S=\sqrt{X^{2}+Y^{2}+Z^{2}}$ 服从 Maxwell 分布律
\begin{align*}p(s)=\sqrt{\frac{2}{\pi}}\frac{s^{2}}{\sigma^{3}}\exp\left(-\frac{s^{2}}{2\sigma^{2}}\right),\qquad s>0.\end{align*}
\end{exercise}

\begin{exercise}\label{exercise:ch03-37}
若 $\xi_1,\xi_2,\cdots,\xi_n$ 相互独立，均服从 $N(0,1)$，试用\eqref{align:ch03-变换密度公式}式，化为 $n$ 维极坐标，证明 $\eta=\xi_1^{2}+\xi_2^{2}+\cdots+\xi_n^{2}$ 服从 $\chi^{2}$ 分布。
\end{exercise}

\begin{exercise}\label{exercise:ch03-38}
若 $\xi$ 与 $\eta$ 相互独立，且分别服从 $\Gamma(r_1,\lambda)$ 及 $\Gamma(r_2,\lambda)$，试求 $\alpha=\xi+\eta$ 与 $\beta=\dfrac{\xi}{\xi+\eta}$ 的联合密度函数 $q(u,v)$，并证明：(1) 随机变量 $\beta$ 服从 $\beta$ 分布：
\begin{align*}p(v)=\frac{\Gamma(r_1+r_2)}{\Gamma(r_1)\,\Gamma(r_2)}v^{r_1-1}(1-v)^{r_2-1},\qquad 0<v<1;\end{align*}
(2) 随机变量 $\alpha$ 与 $\beta$ 独立。
\end{exercise}

\begin{exercise}\label{exercise:ch03-39}
若 $\xi,\eta$ 独立，且均服从 $N(0,1)$，试求 $U=\xi^{2}+\eta^{2}$ 与 $V=\dfrac{\xi}{\eta}$ 的密度函数，并证明它们是独立的。
\end{exercise}

\begin{exercise}\label{exercise:ch03-40}
若 $(\xi,\eta)$ 服从二元正态分布\eqref{align:ch03-二元正态密度}，试找出 $\xi+\eta$ 与 $\xi-\eta$ 相互独立的充要条件。
\end{exercise}

\begin{exercise}\label{exercise:ch03-41}
对二元正态密度函数
\begin{align*}p(x,y)=\frac{1}{2\pi}\exp\left\{-\frac{1}{2}(2x^{2}+y^{2}+2xy-22x-14y+65)\right\}.\end{align*}
(1) 把它化为标准形式\eqref{align:ch03-二元正态密度}；(2) 指出 $\mu_1,\mu_2,\sigma_1,\sigma_2,\rho$；(3) 求 $p_1(x)$；(4) 求 $p(x\mid y)$。
\end{exercise}

\begin{exercise}\label{exercise:ch03-42}
设 $\boldsymbol{\mu}=\boldsymbol{0}$，$\boldsymbol{\Sigma}^{-1}=\begin{pmatrix}7&3&2\\3&4&1\\2&1&2\end{pmatrix}$，试写出分布密度\eqref{align:ch03-多元正态向量}并求出 $(\xi_1,\xi_2)$ 的边际密度函数。
\end{exercise}

\begin{exercise}\label{exercise:ch03-43}
设 $\xi$ 与 $\eta$ 是相互独立相同分布的随机变量，其密度函数不等于 $0$ 且有二阶导数，试证若 $\xi+\eta$ 与 $\xi-\eta$ 相互独立，则随机变量 $\xi,\eta,\xi+\eta,\xi-\eta$ 均服从正态分布。
\end{exercise}

\begin{exercise}\label{exercise:ch03-44}
把习题 25 和习题 26 的结果推广到 $n$ 个随机变量的场合。
\end{exercise}

\begin{exercise}\label{exercise:ch03-45}
设随机变量 $\xi$ 与 $\eta$ 相互独立，$\xi\sim B(n,p)$，而 $\eta$ 服从 $(0,1)$ 上均匀分布，试求 $\xi+\eta$ 的分布函数和密度函数。
\end{exercise}

\begin{exercise}\label{exercise:ch03-46}
试求顺序统计量 $\xi_k^{*}$ 与 $\xi_l^{*}$（$k<l$）的联合密度函数。
\end{exercise}

\begin{exercise}\label{exercise:ch03-47}
试利用概率的连续性重新证明一元分布函数的性质 (ii) 和性质 (iii)，并说明这种证法可推广到多元的场合。
\end{exercise}

\begin{exercise}\label{exercise:ch03-48}
利用随机变量分布解释 Betrand 奇论。
\end{exercise}

\begin{exercise}\label{exercise:ch03-49}
若 $f$ 是 $\varOmega$ 上单值实函数，对 $B\subset\mathbf{R}^{1}$，记 $f^{-1}(B)=\{\omega\in\varOmega:f(\omega)\in B\}$，试证逆映照 $f^{-1}$ 具有如下性质：
\begin{enumerate}
\item $f^{-1}\!\left(\bigcup_{\lambda\in\varLambda}B_\lambda\right)=\bigcup_{\lambda\in\varLambda}f^{-1}(B_\lambda)$；
\item $f^{-1}\!\left(\bigcap_{\lambda\in\varLambda}B_\lambda\right)=\bigcap_{\lambda\in\varLambda}f^{-1}(B_\lambda)$；
\item $f^{-1}(\overline{B})=\overline{f^{-1}(B)}$。
\end{enumerate}
\end{exercise}

\begin{exercise}\label{exercise:ch03-50}
证明：$\xi$ 是一个随机变量当且仅当对任何 $x\in\mathbf{R}^{1}$，成立
\begin{align*}\{\omega:\xi(\omega)<x\}\in\mathscr{F}.\end{align*}

[提示：必要性是显然的，为证充分性，记 $\mathfrak{m}=\{A:A\subset\mathbf{R}^{1},(\xi(\omega)\in A)\in\mathscr{F}\}$，验证 $\mathfrak{m}$ 是 $\sigma$ 域，又 $\mathfrak{m}$ 包含全体形如 $(-\infty,x)$ 的区间，故 $\mathfrak{m}$ 包含 $\mathscr{B}_1$。]
\end{exercise}

\end{document}